% ===================================================================
%  圖表第 8 號 — 收成。
%
%  Traduction, faite en 2026, en cantonais ecrit d'aujourd'hui, en
%  caracteres traditionnels, sur l'ido de texto/io. Regles de la
%  colonne : voir 10-toubiu-01.tex. Deux scenes, deux numerotations :
%  les cles portent c1 et c2.
%
%  DEUX BETES DONNENT ICI LEUR MOT CANTONAIS, ET AUCUN DES DEUX NE
%  S'ECRIT AU NORD :
%
%    蠄蟝 (1)   le crapaud   Pekin : 蟾蜍
%    蛤乸 (17)  la grenouille Pekin : 青蛙
%
%  « 蠄蟝 » est un mot a deux caracteres qui ne servent qu'a lui ;
%  « 蛤乸 » reprend le suffixe feminin 乸 releve au tableau 7, ou il
%  faisait 雞乸 et 豬乸 — la grenouille est nommee au feminin comme
%  la poule, sans que personne pense au sexe de la bete.
%
%  ET LE CERISIER PORTE LE MOT LE PLUS VOYANT DU TABLEAU :
%
%    車厘子 (5)  la cerise   Pekin : 樱桃
%    車厘子樹 (4) le cerisier Pekin : 樱桃树
%
%  C'est l'anglais « cherries » entre par Hong Kong, decoupe en trois
%  syllabes cantonaises. Le meme procede que 芝士 au tableau 7 : la
%  ou Pekin traduit, Canton transcrit. Le fruit du marche s'appelle
%  ainsi et pas autrement.
%
%  LES METIERS CONTINUENT EN 佬 : 收割佬 (26), 箍桶佬 (19),
%  摘葡萄佬 (18), 打獵佬 (5), 釣魚佬 (44), 撒網佬 (48). Le tableau
%  en aligne six de plus, apres les six du tableau 7 et les neuf du
%  batiment au tableau 6.
%
%  L'EMBARQUEMENT SE DIT 艇 ET NON 船 (49) ; la bicyclette 單車
%  comme au tableau 3 ; le ruisseau 河仔 comme au tableau 7.
%
%  TROIS BLOCS ONT DEMANDE L'APPOSITION OU L'INVERSION PREVENTIVE :
%  c1-08-1 (les culbutes AVANT les tas de paille, et la faucille
%  avant sa pierre a aiguiser), c1-10-1 (le seigle avant la
%  moissonneuse) et c1-11-1 (la batteuse avant sa locomobile). Le
%  chinois aurait mis chaque fois l'instrument devant l'ouvrage.
%
%  ECART DECLARE : le bloc c2-03-1 rend le renvoi (13) du gobelet,
%  que le fac-simile ido compose « 13) » sans parenthese ouvrante.
%  renvoji.py le passe, la raison est inscrite dans APARTA.
% ===================================================================

%%K t08-noto-1 noto
\VUnotes{143.2mm}{%
\textit{(*) 因為呢個字唔夠準：究竟係收亞麻、收葡萄，定係收蘋果呢？}\nl
\textit{或者用} \VUgras{mesono} \textit{會好啲，呢個字喺《Mondo》第十四期}\nl
\textit{242 頁提出過。不過到而家為止，呢個新詞根仲未畀學院採納，}\nl
\textit{仲要照伊多語慣常嘅規矩等人討論。}%
}

%%K t08-apar-1 apar
\VUsaut{5.0mm}
\VUtitre{15.0pt}{60}{圖表第 8 號}
\VUsaut{3.0mm}
\VUcentre{13.2pt}{90}{\textit{第一場。}}
\VUsaut{3.0mm}
\VUpk{10.2pt}{40}{收成 (*)。}
\VUsaut{4.0mm}
\VUpk{10.2pt}{40}{郊野嘅景色。}
\VUsaut{3.0mm}

%%K t08-c1-01-1 p
1. --- \VUgras{夏天}一到，郊野嘅樣又變過。落咗喺犁溝度嘅種子發咗芽；青青哋
嘅細苗鑽出泥面，跟住抽穗。穀粒結咗出嚟，之後熟晒，而家就淨係差收割。

%%K t08-c1-02-1 p
2. --- \VUgras{野花}開晒。\VUgras{矢車菊} \textsuperscript{(1)}、
\VUgras{虞美人} \textsuperscript{(2)} 同\VUgras{毛地黃} \textsuperscript{(3)} 喺一片
一式一樣嘅麥田色度，加返幾筆鮮嫩\VUgras{花冠}嘅開懷顏色。

%%K t08-c1-03-1 p
3. --- 啲果樹長晒葉，啲果都熟嘞。細細棵\VUgras{車厘子樹} \textsuperscript{(4)} 結到
成樹靚\VUgras{車厘子} \textsuperscript{(5)}，啲細路摘得又食得，開心到不得了。

%%K t08-c1-04-1 p
4. --- 由今日起，啲牲口成日都可以留喺露天度。

%%K t08-c1-04-2 p
啲\VUgras{羊仔} \textsuperscript{(6)} 大咗，變成靚\VUgras{羊乸} \textsuperscript{(7)} 同靚
\VUgras{羊公} \textsuperscript{(8)}，一身毛厚笠笠。佢哋靜靜哋喺\VUgras{牧場}
\textsuperscript{(9)} 度食\VUgras{草，}成塊地畀\VUgras{雛菊}點到花花斑斑。

%%K t08-c1-04-3 p
好快就要幫\VUgras{呢啲}牲口剪毛嘞，攞佢哋啲\VUgras{羊毛}去織我哋着嗰啲布。

%%K t08-tit-2 sub
\VUsaut{5.0mm}
\VUpk{10.2pt}{40}{牧羊人。}
\VUsaut{3.4mm}

%%K t08-c1-05-1 p
5. --- \VUgras{牧羊人} \textsuperscript{(10)} 好似唔點理佢哋。佢淨係望住天邊掠過
嗰場風暴。至於個\VUgras{睇羊嘅} \textsuperscript{(13)}，佢揞住個\VUgras{水壺}
\textsuperscript{(14)} 埋口飲，睇落仲要無所謂。

%%K t08-c1-06-1 p
6. --- 熱到侷住；連隻\VUgras{狗} \textsuperscript{(15)} 都行埋去解渴，佢
𠹌住\VUgras{河仔} \textsuperscript{(16)} 嘅涼水，嚇親隻匿埋喺岸邊草叢嘅可憐\VUgras{蛤乸}
\textsuperscript{(17)}。隻蛤乸一嘢彈落清水度，水面開住啲\VUgras{睡蓮}
\textsuperscript{(18)}。

%%K t08-c1-07-1 p
7. --- 一隻\VUgras{鶺鴒} \textsuperscript{(19)} 喺\VUgras{泉眼} \textsuperscript{(20)}
邊沖涼，度泉水喺兩嚿石中間噴出嚟；佢仲匿埋喺\VUgras{有刺嘅矮樹}
\textsuperscript{(21)} 同啲\VUgras{山楂叢} \textsuperscript{(22)} 下面。

%%K t08-tit-3 sub
\VUsaut{5.0mm}
\VUpk{10.2pt}{40}{收成。}
\VUsaut{3.4mm}

%%K t08-c1-08-1 p
8. --- 對鄉下細路嚟講，割麥係一年之中最好玩嘅日子。佢哋開開心心噉
\VUgras{打筋斗}\textsuperscript{(24)}，就喺啲\VUgras{禾稈堆} \textsuperscript{(23)} 度；
佢哋又幫手\VUgras{收割佬} \textsuperscript{(26)}。呢啲人一路割緊塊\VUgras{麥田}
\textsuperscript{(27)}，佢哋把\VUgras{鐮刀} \textsuperscript{(25)} 畀\VUgras{磨刀石}
\textsuperscript{(30)} 磨到好鋒利；啲細路就執麥，紮成\VUgras{禾把}
\textsuperscript{(31)} 或者\VUgras{細紮} \textsuperscript{(32)}。

%%K t08-c1-09-1 p
9. --- 有時大人會畀最大嗰個攞把\VUgras{鐮仔} \textsuperscript{(12)} 割啲
\VUgras{金黃嘅稈} \textsuperscript{(28)}，稈上面孭住沉甸甸、滿係穀粒嘅\VUgras{麥穗}
\textsuperscript{(29)}。細嗰啲就睇住啲\VUgras{牲口} \textsuperscript{(33)}，佢哋喺
\VUgras{牧場} \textsuperscript{(34)} 度食草，就喺棵大\VUgras{椴樹} \textsuperscript{(35)}
嘅樹蔭下面；細路就用個\VUgras{網仔} \textsuperscript{(36)} 捉靚\VUgras{蝴蝶。}

%%K t08-c1-09-2 p
仲有幾個爬上架\VUgras{車仔} \textsuperscript{(37)} 度，人哋用\VUgras{草叉}
\textsuperscript{(38)} 遞禾把上嚟，佢哋就負責擺好。

%%K t08-c1-10-1 p
10. --- 成片\VUgras{平原} \textsuperscript{(39)} 都熱鬧晒，啲\VUgras{雲雀}
\textsuperscript{(41)} 喺上面飛起飛落。個個都忙住收成。不過窮嗰啲仲用緊舊家生
——\VUgras{鐮刀、鐮仔}同\VUgras{連枷} \textsuperscript{(40)}；有錢嘅地主就叫人割佢哋
啲麥同\VUgras{黑麥} \textsuperscript{(43)}，用嘅係\VUgras{收割機} \textsuperscript{(42)}。
之後唔使將禾把運返穀倉，就地砌成一個個大\VUgras{禾堆} \textsuperscript{(44)}。

%%K t08-c1-11-1 p
11. --- 跟住就拉架\VUgras{打穀機} \textsuperscript{(47)} 埋嚟，架
\VUgras{蒸汽機} \textsuperscript{(45)} 用條\VUgras{傳動皮帶} \textsuperscript{(46)} 帶動佢；
噉樣穀就同\VUgras{禾稈}分開咗，都唔使離開播佢落去嗰塊田。

%%K t08-tit-4 sub
\VUsaut{5.0mm}
\VUpk{10.2pt}{40}{風暴。}
\VUsaut{3.4mm}

%%K t08-c1-12-1 p
12. --- 收成嗰陣時，成日有\VUgras{大風暴} \textsuperscript{(53)}。首先聽到遠處
\VUgras{行雷}隆隆聲；跟住\VUgras{西邊嘅大風} \textsuperscript{(54)} 吹起，呼呼哋
吹到\VUgras{樹林} \textsuperscript{(49)} 入面最粗嗰啲樹都彎晒腰。黑色嘅\VUgras{烏雲}
\textsuperscript{(56)} 攻上天，一道道刺眼嘅\VUgras{閃電} \textsuperscript{(55)} 打橫劃過。
\VUgras{雨}落嘞，有時仲落\VUgras{雹，}將就快收得嘅莊稼壓到扁晒。

%%K t08-c1-13-1 p
13. --- 閃電時時燒着屋。舊年就係噉，\VUgras{風車} \textsuperscript{(50)} 個頂燒到
淨返灰，佢兩塊\VUgras{風翼} \textsuperscript{(51)} 都撞爛咗。

%%K t08-c1-14-1 p
14. --- 過幾個鐘，天又靜返。天邊出咗條光閃閃嘅\VUgras{彩虹} \textsuperscript{(52)}，
大自然又接返啱啱斷咗一陣嘅生活。啲鄉下人本來匿咗入啲\VUgras{茅屋}
\textsuperscript{(48)}，而家又出返嚟做嘢，好硬淨噉去執返頭先損失嗰啲。

%%K t08-tit-5 sub
\VUsaut{6.0mm}
\VUcentre{13.2pt}{90}{\textit{第二場。}}
\VUsaut{3.6mm}

%%K t08-tit-6 sub
\VUpk{10.2pt}{40}{打獵。--- 摘葡萄。}
\VUsaut{1.4mm}
\VUpk{10.2pt}{40}{釣魚。}
\VUsaut{4.0mm}

%%K t08-c2-01-1 p
1. --- 大約\VUgras{八月}尾或者\VUgras{九月}中，睇地方而定，就開獵期嘞。第一道
日頭先啱啱令到啲\VUgras{蜥蜴} \textsuperscript{(2)} 出返個窿口，個\VUgras{打獵佬}
\textsuperscript{(5)} 就已經帶住佢啲\VUgras{狗} \textsuperscript{(4)} 上路。

%%K t08-c2-02-1 p
2. --- 佢着咗套厚布做嘅結實\VUgras{獵裝} \textsuperscript{(9)}，腳上綁住皮
\VUgras{綁腿，}噉先至行得入啲濕草度——啲\VUgras{蠄蟝} \textsuperscript{(1)} 就匿喺嗰度；
佢攞埋支\VUgras{獵槍} \textsuperscript{(6)} 同個\VUgras{獵物袋} \textsuperscript{(10)}，就出發嘞。

%%K t08-c2-03-1 p
3. --- 佢一心追獵物，就走到葡萄園中間，嗰度摘葡萄摘得好熱鬧。種葡萄佬啲
細路平時要喺路邊睇住啲山羊，今日就由得佢哋周圍食草；佢哋將隻老
\VUgras{山羊公} \textsuperscript{(3)} 綁咗喺條柱度——嗰隻梗係唔會放過機會走人嘅——
之後就自己攞返一份熱鬧。有個喺個\VUgras{摘葡萄籃} \textsuperscript{(12)} 度捉住一梳葡萄，
玩住將啲\VUgras{汁} \textsuperscript{(14)} 倒落個細\VUgras{杯} \textsuperscript{(13)} 度，扮做
葡萄汁（即係未發酵嘅酒）。另一個就將啱啱織好嘅\VUgras{葉冠} \textsuperscript{(15)}
戴上頭。

%%K t08-c2-03-1 p suite
佢哋隔籬，一啲唔怕醜嘅\VUgras{麻雀} \textsuperscript{(16)} 走到榨汁機前面偷葡萄食。

%%K t08-c2-04-1 p
4. --- 細路仔玩嘅時候，大人就做嘢。\VUgras{摘葡萄佬} \textsuperscript{(18)} 用
\VUgras{揹籮} \textsuperscript{(17)} 將葡萄擔入酒窖；\VUgras{箍桶佬} \textsuperscript{(19)}
喺度整啲\VUgras{木桶} \textsuperscript{(20)}，睇下啲\VUgras{桶板} \textsuperscript{(22)} 同啲
\VUgras{桶底} \textsuperscript{(21)} 妥唔妥當，又換走爛咗嘅\VUgras{鐵箍}
\textsuperscript{(23)}。啲\VUgras{大木盆} \textsuperscript{(25)} 都係佢做嘅，人哋將啲
\VUgras{葡萄} \textsuperscript{(26)} 倒落去發酵。

%%K t08-c2-05-1 p
5. --- 發酵完，就抽走啲酒，將啲渣擺上\VUgras{榨汁機} \textsuperscript{(28)}，慢慢
扭緊條\VUgras{大螺絲} \textsuperscript{(29)} 榨出啲汁，汁就流落個\VUgras{細盆}
\textsuperscript{(32)} 度，仲可以再攞多啲\VUgras{酒} \textsuperscript{(31)}。

%%K t08-tit-8 sub
\VUsaut{6.0mm}
\VUpk{10.2pt}{40}{啤酒同蘋果酒。}
\VUsaut{3.4mm}

%%K t08-c2-06-1 p
6. --- \VUgras{酒窖} \textsuperscript{(27)} 幅牆上面爬住一枝\VUgras{啤酒花}
\textsuperscript{(30)}。喺呢度佢淨係做裝飾，不過有啲地方好少葡萄，就特登種啤酒花
嚟釀\VUgras{啤酒。}

%%K t08-c2-07-1 p
7. --- 細細棵\VUgras{蘋果樹} \textsuperscript{(33)} 結到成樹蘋果，熟咗就可以釀出
一流嘅\VUgras{蘋果酒。}喺個\VUgras{籠} \textsuperscript{(34)} 入面，隻\VUgras{金絲雀}
\textsuperscript{(35)} 同隻\VUgras{金翅雀} \textsuperscript{(36)} 眼紅紅噉望住啲自出自入
嘅麻雀。

%%K t08-c2-08-1 p
8. --- 一路望到天邊，山坡上面排住一行行\VUgras{葡萄架} \textsuperscript{(40)}，撐住
啲靚到極嘅\VUgras{葡萄藤} \textsuperscript{(39)}，藤上掛滿沉甸甸嘅\VUgras{葡萄梳。}

%%K t08-c2-08-2 p
成年落嚟，\VUgras{種葡萄佬} \textsuperscript{(42)} 都喺度打理佢個\VUgras{葡萄園}
\textsuperscript{(38)}，而家就得返一造靚收成做報酬。啲\VUgras{摘葡萄婆}
\textsuperscript{(41)} 將車上面啲盆裝到滿瀉，架車就由啲\VUgras{騾} \textsuperscript{(43)}
拉住，個個都笑騎騎。

%%K t08-tit-9 sub
\VUsaut{5.0mm}
\VUpk{10.2pt}{40}{釣魚。}
\VUsaut{3.4mm}

%%K t08-c2-09-1 p
9. --- 啲\VUgras{釣魚佬} \textsuperscript{(44)} 都係一樣，佢哋一早就攞住啲
\VUgras{魚竿} \textsuperscript{(45)} 嚟到水邊坐低。

%%K t08-c2-09-2 p
啲\VUgras{魚} \textsuperscript{(47)} 一見到佢哋將條\VUgras{魚絲} \textsuperscript{(46)}
放落水就咬\VUgras{魚鈎；}佢哋仲未嚟得切換過條\VUgras{魚餌，}又有新一條
上釣，喺魚絲尾度彈跳。

%%K t08-c2-09-3 p
不過\VUgras{撒網佬} \textsuperscript{(48)} 就捉得多啲：佢企喺自己隻\VUgras{艇}
\textsuperscript{(49)} 度，將個\VUgras{網}撒落\VUgras{條河} \textsuperscript{(50)} 中間。

%%K t08-c2-10-1 p
10. --- 秋天係最啱去散步嘅季節：行路、\VUgras{騎馬} \textsuperscript{(52)}、踩
\VUgras{單車} \textsuperscript{(53)}、揸\VUgras{汽車} \textsuperscript{(54)} 都得。啲
\VUgras{路} \textsuperscript{(51)} 乾乾淨淨，大家就趁最後幾日好天。因為啲\VUgras{燕子}
\textsuperscript{(56)} 已經喺屋頂度埋堆，就快拍晒對翼飛去暖啲嘅地方。老棵
\VUgras{核桃樹} \textsuperscript{(57)} 啲葉黃咗，啲\VUgras{核桃}跌落地；啲高
\VUgras{樹}一叢叢好似一束\VUgras{花} \textsuperscript{(58)} 噉企喺\VUgras{高地}
\textsuperscript{(59)} 上面，都快啲落清啲葉嘞。

%%K t08-c2-11-1 p
11. --- \VUgras{太陽} \textsuperscript{(60)} 遲咗先出，日頭短咗，一早
\VUgras{朝頭早}就有股相當刺骨嘅凍意；而家一群群\VUgras{山鷸} \textsuperscript{(61)}
已經離開我哋呢啲留唔住客嘅天氣嘞。

\VUsaut{6.0mm}
\VUfilet{22mm}
